%! Piecewise-Defined Functions — Unit 2, Lesson 2.6 — Student Packet Cover Sheet
\documentclass[10pt]{article}
\usepackage{atda-article}
\usepackage{atda-boxes}
\usepackage{ltablex}
\keepXColumns

\begin{document}

% Full-bleed royal banner
\begin{tikzpicture}[remember picture, overlay]
  \fill[royal] (current page.north west)
    rectangle ([yshift=-0.9in]current page.north east);
\end{tikzpicture}
\vspace{-0.8in}

\begin{center}
  {\color{white}\textbf{\LARGE Algebra, Trigonometry, and Data Analysis}}\\[4pt]
  {\color{mistmid}\large\textbf{Unit 2: Functions, Transformations, and Graph Analysis}}\\[6pt]
  {\color{white}\textbf{\large Lesson 2.6 \quad Piecewise-Defined Functions}}
\end{center}

\vspace{0.22in}
\namedateperiod
\vspace{0.18in}

\begin{learningtargetbox}
\small
\begin{enumerate}[leftmargin=1.5em, itemsep=5pt]
  \item \ldots{} read a graph that takes \textbf{more than one rule}, and evaluate it by first asking
        \emph{which piece owns this input} and then using that piece's rule and only that one.
  \item \ldots{} explain why a piecewise-defined function is still \textbf{one} function: every input
        in the domain belongs to exactly one piece, so every input still has exactly one output.
  \item \ldots{} use the \textbf{open and closed circles} at a boundary to say what the function's
        value is there --- and use them to write a range with a bracket or a parenthesis.
  \item \ldots{} decide whether a graph is \textbf{continuous} at a boundary --- can it be traced
        without lifting the pencil? --- and tell a \emph{corner} from a \emph{break}.
  \item \ldots{} find all of Unit 2's characteristics on a piecewise graph: domain and range, zeros
        and intercepts, increasing/decreasing/constant intervals, absolute maximum and minimum, end
        behavior, and asymptotes.
\end{enumerate}
\end{learningtargetbox}

\vspace{0.14in}

\begin{tocbox}
\small
\renewcommand{\arraystretch}{1.5}
\begin{tabularx}{\linewidth}{@{} c l X r @{}}
\rowcolor{mist}
\textbf{\#} & \textbf{Component} & \textbf{Description} & \textbf{Score} \\
\midrule
1 & Warm-Up        & Two rules, two stretches, and a graph you can trace without lifting your pencil
                                                                   & \blank{1.2cm} \\
2 & Guided Notes   & What a piecewise function is, where each rule lives, and what the dots at a
                     boundary are telling you                       & \blank{1.2cm} \\
3 & Group Activity & A phone plan, a parking garage priced in steps, and a graph that jumps
                                                                   & \blank{1.2cm} \\
4 & Exit Ticket    & One babysitting rate, and a classmate who called a corner a break
                                                                   & \blank{1.2cm} \\
5 & Homework       & A tiered electricity bill, a range that is a list, and a look ahead
                                                                   & \blank{1.2cm} \\
\midrule
  & \multicolumn{2}{r}{\textbf{Total}} & \blank{1.2cm} \\
\end{tabularx}
\end{tocbox}

\vspace{0.14in}

\begin{remindbox}
\small
This lesson covers \texttt{AFDA.AF.2} applied to \textbf{piecewise-defined} functions --- every
characteristic you have met in this unit, now read off a graph that takes more than one rule.
Everything is still read \textbf{off a graph}. One sentence runs the whole packet: \textbf{a piecewise
function is one function, not several.} So before you use a rule, ask which stretch of the domain it
owns --- and when you reach a boundary, \textbf{let the dots decide}: a filled circle is included, an
open circle is not.
\end{remindbox}

\end{document}
